\chapter{API Design and Implementation}

\section{Introduction}

All functionality of OrchardVision is exposed through a REST API under the \code{/api/v1/} prefix. The same API is used by the React client and can be called directly by scripts. This chapter lists the endpoints, explains the request flows, and describes error handling and the shared TypeScript types.

\section{Endpoint Summary}

Table~\ref{tab:api} summarises the endpoints.

{\singlespacing\small
\begin{longtable}{@{}p{1.5cm}p{6.6cm}p{5.8cm}@{}}
\caption{REST API endpoints.}\label{tab:api}\\
\toprule
\textbf{Method} & \textbf{Path} & \textbf{Description}\\
\midrule
\endfirsthead
\toprule
\textbf{Method} & \textbf{Path} & \textbf{Description}\\
\midrule
\endhead
\bottomrule
\endfoot
GET & \code{/health} & Service health\\
POST & \code{/videos} & Upload video\\
GET & \code{/videos} & List videos\\
GET & \code{/videos/{id}} & Video metadata\\
GET & \code{/videos/{id}/tracks} & Track trajectories\\
GET & \code{/videos/{id}/identities} & Identity prototypes\\
GET & \code{/videos/{id}/identities/projection} & 2D coordinates for visualisation\\
GET & \code{/videos/{id}/identities/{identity_id}} & Specific identity\\
GET & \code{/videos/{id}/identities/{identity_id}/observations} & Per-frame observations\\
GET & \code{/videos/{id}/reid-events} & Re-identification events\\
GET & \code{/videos/{id}/detections} & Persisted detections\\
GET & \code{/videos/{id}/per-frame} & Per-frame timing telemetry\\
GET & \code{/videos/{id}/motion} & Camera motion estimate\\
GET & \code{/videos/{id}/analytics} & Aggregated analytics\\
DELETE & \code{/videos/{id}} & Delete video and all associated data\\
POST & \code{/videos/{id}/process} & Start processing job\\
GET & \code{/jobs/{id}} & Job status\\
GET & \code{/experiments} & List experiments\\
POST & \code{/experiments} & Create and run experiment\\
GET & \code{/experiments/{id}} & Experiment detail\\
GET & \code{/experiments/{id}/export} & Download results (JSON/CSV)\\
DELETE & \code{/experiments/{id}} & Delete experiment (cannot be \code{RUNNING})\\
GET & \code{/settings} & Pipeline defaults\\
PUT & \code{/settings} & Update pipeline defaults\\
\end{longtable}}

\section{Resource-to-Model Mapping}

Each read endpoint returns data that the pipeline persisted in one of the models of Chapter~4 (Table~\ref{tab:map}).

\begin{table}[htbp]
\centering\small
\begin{tabular}{@{}>{\raggedright\arraybackslash}p{5.9cm}>{\raggedright\arraybackslash}p{3.7cm}>{\raggedright\arraybackslash}p{3.8cm}@{}}
\toprule
\textbf{Endpoint} & \textbf{Model} & \textbf{Created in stage}\\
\midrule
\code{/videos/{id}} & \code{Video} & Upload\\
\code{/videos/{id}/tracks} & \code{Track} & Tracking\\
\code{/videos/{id}/detections} & \code{DetectionRecord} & Detection\\
\code{/videos/{id}/identities} & \code{IdentityRecord} & Identity consolidation\\
\code{/videos/{id}/identities/{identity_id}/observations} & \code{IdentityObservation} & Identity consolidation\\
\code{/videos/{id}/reid-events} & \code{ReidEvent} & Re-identification\\
\code{/videos/{id}/per-frame} & \code{FrameTimingRecord} & All stages\\
\code{/videos/{id}/motion} & \code{CameraMotionRecord} & Camera-motion estimation\\
\code{/videos/{id}/analytics} & Aggregated metrics & Analytics\\
\code{/jobs/{id}} & \code{ProcessingJob} & Process request\\
\code{/experiments} & \code{Experiment} & Experiment request\\
\code{/settings} & \code{AppSettings} & Settings update\\
\bottomrule
\end{tabular}
\caption[Endpoint-to-model mapping]{Mapping between endpoints and the persisted models. \add{The ``stage'' column is inferred from the pipeline description.}}
\label{tab:map}
\end{table}

\figph{fig:seq}{Sequence diagram of a processing request}{Sequence of calls for uploading and processing a video.}{UML sequence diagram with lifelines: Client, FastAPI router, StorageService, Database, Celery/background task, ProcessingPipeline. Messages: POST /videos $\rightarrow$ save\_upload $\rightarrow$ extract\_video\_metadata $\rightarrow$ insert Video (uploaded); POST /videos/\{id\}/process $\rightarrow$ validate no active job $\rightarrow$ insert ProcessingJob (QUEUED) $\rightarrow$ enqueue; worker updates stage and progress; final persistence and Video status processed; client polls GET /jobs/\{id\}.}{7cm}

\section{Request Flows}

\subsection{Video Upload}

\code{POST /videos} saves the upload through \code{storage_service.save_upload}, extracts metadata with \code{extract_video_metadata} and creates the \code{Video} row with status \code{uploaded}. The maximum upload size is 2\,GB by default and can be changed through \code{MAX_UPLOAD_BYTES}.

\subsection{Starting a Processing Job}

\code{POST /videos/{id}/process} validates that no job is active for the video, creates a \code{ProcessingJob} with status \code{QUEUED}, and then enqueues it with Celery or runs it synchronously as a background task. The client polls \code{GET /jobs/{id}} to read the job's \code{status}, \code{stage} and \code{progress}.

\subsection{Reading Results}

After a video is \code{processed}, its tracks, identities, observations, ReID events, detections, per-frame timing, camera motion and analytics are retrieved with the GET endpoints in Table~\ref{tab:api}. The identity projection endpoint returns two-dimensional coordinates computed by the \code{EmbeddingProjection} service for plotting.

\subsection{Experiments}

\code{POST /experiments} creates and runs an experiment described by a name, a tracker, a scenario and a number of frames. Its results can be read with \code{GET /experiments/{id}} and downloaded as JSON or CSV with \code{GET /experiments/{id}/export}. An experiment that is \code{RUNNING} cannot be deleted.

\subsection{Settings}

\code{GET /settings} and \code{PUT /settings} read and update the pipeline defaults stored in the single \code{AppSettings} row.

\section{Error Handling}

All errors are returned in a single format:

\begin{lstlisting}[style=plain]
{"detail": "...", "code": "...", "status": ...}
\end{lstlisting}

The API defines a base \code{ApiError} and the custom exceptions \code{NotFoundError}, \code{ConflictError} and \code{ValidationFailError}. A global exception handler in \code{apps/api/app/main.py} converts them into the response format above.

\section{Shared TypeScript Types}

Common interfaces are defined in \code{packages/types/src/common.ts} and are used by the frontend and by any other TypeScript client:

\begin{itemize}
  \item \code{BoundingBox}: \code{{ x: number, y: number, width: number, height: number }};
  \item \code{VisibilityState}: \code{'visible' | 'occluded' | 'lost'};
  \item \code{JobStatus}: \code{'uploaded' | 'processing' | 'processed' | 'error'};
  \item \code{ProcessingStage}: \code{'frame_extraction'}, \code{'detection'}, \code{'tracking'}, \code{'reid'}, \code{'analytics'};
  \item \code{ProcessingStatus}: \code{'QUEUED'}, \code{'RUNNING'}, \code{'COMPLETED'}, \code{'FAILED'};
  \item response wrappers such as \code{MultiVideoResponse}, \code{MultiTrackResponse} and \code{MultiIdentityResponse}.
\end{itemize}

The types can be checked with \code{pnpm exec tsc --noEmit} from \code{packages/types/}.

\section{Frontend Integration}

The components in \code{apps/web/src/} use an API client that calls the FastAPI backend. During development the Vite dev server (\code{npm run dev}) proxies requests under \code{/api/} to the backend, and the backend's CORS configuration allows \code{localhost:5173}. \add{The interface is organised around the main user tasks: uploading and listing videos, starting and monitoring processing, inspecting tracks and identities, viewing the identity projection, managing experiments and editing the pipeline settings.}

\section{Example Interaction}

\add{The following listing shows a typical sequence of calls for processing a video. The identifiers and the request body are illustrative.}

\begin{lstlisting}[style=added,language=bash]
# 1. upload a video
curl -F "file=@row_03.mp4" http://localhost:8000/api/v1/videos
# 2. start processing video 1 (tracker chosen in the body)
curl -X POST http://localhost:8000/api/v1/videos/1/process \
     -H "Content-Type: application/json" \
     -d '{"tracker": "bytetrack"}'
# 3. poll the job
curl http://localhost:8000/api/v1/jobs/1
# 4. read results
curl http://localhost:8000/api/v1/videos/1/analytics
\end{lstlisting}
